\begin{statement}{C.25}{Unrestricted cusp conjugation changes the D.45 finite test}{Refuted}
The assertion that the recorded sphere action of LM7 D.45
has no peripheral lift from
\(\mathcal B/D_9(F)\) to \(\mathcal B/D_{10}(F)\)
is false if ``peripheral'' allows an arbitrary conjugate
of the cusp. There is an actual automorphism of the latter
finite group inducing exactly that lower action and sending
\(s,r,\tau\) to conjugates of themselves.

The strongest exclusion established by that test fixes
the cusp frame to one: with this normalization, all
torsion-frame and compatible multiplier changes fail,
already on the explicit central \(C_2\)-quotient in
C.12. The finite counterexample does not decide
existence on the whole relative completion.
\end{statement}
\begin{proof}
Use the literal words \(b,a,v\) and the automorphism
\(\widetilde\alpha\) of (C.12h)--(C.12i).
The independently run program
\path{c_cusp_centralizer_scope.py} verifies the complete
Magnus equations, the unchanged lower action, and the
Frattini identity proving surjectivity. The torsion
images are conjugates by construction, and
\(\widetilde\alpha(\tau)=v\tau v^{-1}\).
The lower cusp-centralizer's 23 generators enlarge
the correction rank from 47 to 56; the saved solution
is an explicit witness, not only a dimension count.
C.12 proves the normalized exclusion and constructs
its central quotient. Both claims therefore concern
precisely specified, different finite lifting problems.
\end{proof}

\begin{statement}{C.26}{The folded four-strand pentagon test is identically blind}{Proved}
For every \(f\in\widehat F_2'\) satisfying the two-cycle
relation \(f(y,x)f(x,y)=1\), the four-strand pentagon
holds after the specialization \(\sigma_1=\sigma_3=T\),
\(\sigma_2=U\) of (C.26a), at every finite congruence
modulus. In particular this test cannot force
\(\Delta=0\) on the gentle group.
\end{statement}
\begin{proof}
Take the actual braid homomorphism
\[
 B_4\longrightarrow{\rm SL}_2(\mathbb Z),\qquad
 \sigma_1,\sigma_3\longmapsto T=
       \begin{pmatrix}1&1\\0&1\end{pmatrix},\qquad
 \sigma_2\longmapsto U=
       \begin{pmatrix}1&0\\-1&1\end{pmatrix}.
 \tag{C.26a}
\]
Set \(a=T^2,b=U^2,c=UT^2U^{-1},d=TU^2T^{-1}\),
the respective images of \(x_{12},x_{23},x_{13},x_{24}\);
\(x_{34}\) has image \(a\). Exact integral multiplication gives
\[
 ac=da=-b^{-1},\qquad bd=cb=-a^{-1}.
 \tag{C.26b}
\]
The five terms of the profinite pentagon
\citep[Introduction, displayed pentagon]{cDolgushevGuynee2024}
therefore specialize to
\[
 f(b,a),\quad f(-b^{-1},-b^{-1}),\quad f(a,b),\quad
 f(a,-a^{-1}),\quad f(-a^{-1},a).
 \tag{C.26c}
\]
The three terms with commuting arguments are one because
\(f\in\widehat F_2'\), and the remaining two cancel by
the gentle two-cycle relation. Thus this specialization
is satisfied by \emph{every} gentle pair at every modulus,
including the E.58 pair with \(\Delta=1\).
The exact integer certificate \path{c_folded_pentagon.py}
was run. This proves that the folded two-dimensional
four-strand test cannot decide the remaining parity
question; no such conclusion is drawn for other
four-strand quotients.

\end{proof}

\begin{statement}{C.27}{Exact affine pentagon sieves on common word images}{Proved}
Consider the distinct-endpoint braid representation into
\((\mathbb Z/m)^2\rtimes{\rm SL}_2(\mathbb Z/m)\)
given by
\[
 \sigma_1\mapsto(0,T),\qquad
 \sigma_2\mapsto(0,U),\qquad
 \sigma_3\mapsto(e_1,T),
 \tag{C.27a}
\]
with \(T,U\) as in (C.26a). Use all five pentagon
evaluation pairs, and the necessary multiplier-one
two- and three-cycle identities at their six
permuted triples. Let \(D_m\) be the actual common
image of \(\widehat F_2'\), let \(C_m\subseteq D_m\)
satisfy those necessary gentle identities, and let
\(P_m\subseteq C_m\) also satisfy the pentagon.
The following exact counts hold:
\[
\begin{array}{c|r|r|r|r}
 m&|D_m|&|C_m|&|P_m|&
   |\operatorname{ev}_{(x_{12},x_{23})}(C_m)|\\\hline
 3&216&9&3&3\\
 4&1&1&1&1\\
 5&15000&25&5&5\\
 8&2&1&1&1
\end{array}
\tag{C.27b}
\]
In each row the base-evaluation image of \(P_m\)
equals that of \(C_m\). These are exhaustive common
word-image tests, not independently chosen values
of the five factors. They do not assert that every
member of \(C_m\) is a full gentle-group image.
\end{statement}
\begin{proof}
Direct semidirect-product multiplication proves both
adjacent braid relations and the far-commutation relation
in (C.27a). Put
\(a=\sigma_1^2,b=\sigma_2^2,e=\sigma_3^2\),
\(c=\sigma_2a\sigma_2^{-1}\),
\(d=\sigma_3b\sigma_3^{-1}\).
The five pairs are
\[
 (b,e),\quad(ac,de),\quad(a,b),\quad(a,bd),\quad(cb,e).
\]
For each pair \((x,y)\), put \(z=(xy)^{-1}\),
and include the six pairs
\[
 (x,y),\ (y,x),\ (z,x),\ (x,z),\ (y,z),\ (z,y).
\]
At each of them the program includes the additional
evaluations needed to test
\[
 f(y,x)f(x,y)=1,
 \qquad f(z,x)f(y,z)f(x,y)=1.
\]
Every genuine multiplier-one gentle pair satisfies
these equations, so they are necessary conditions
with their exact scope explicit.

The resulting two diagonal source generators are
\(X=(x_j)_j\) and \(Y=(y_j)_j\).
The first program, \path{c_affine_pentagon_probe.py},
constructs their faithful block permutation action,
computes \(\langle X,Y\rangle'\), and enumerates
every element before applying the equations.
The independent program
\path{c_affine_pentagon_verify.py} imports no group
theory package. It enumerates the normal closure of
\(C=[X,Y]\) using the three maps
\[
 h\longmapsto hC,\qquad
 h\longmapsto XhX^{-1},\qquad
 h\longmapsto YhY^{-1}.
\]
On the finite ambient group these maps are permutations.
Their forward orbit of one is therefore closed under
their inverses; it is exactly the normal closure of
\(C\), hence the whole derived image.
Every discovered element retains an actual commutator
word. The program evaluates the literal words for
every survivor and every excluded necessary-gentle
candidate, verifies zero abelianization, and recovers
all counts in (C.27b), including equality of the two
base-image sets. Complete words and values are saved
in the corresponding JSON certificates.

The modulus-three representation is the Ferrari
homology representation already used in
\emph{LM7\_PUSH E.17}; this calculation independently
checks the stated finite sieve and does not claim a new
whole-CM-fibre theorem. In particular the base value
\(f(a,b)\) alone does not determine the normalized
CM frames. The finite counts supply no answer to the
remaining full-GT parity question in C.15.
\end{proof}

\begin{statement}{C.28}{Perfect common braid images cannot exclude solvable shadows}{Proved}
Let \(\eta:F_2\twoheadrightarrow S\) be any finite
solvable quotient, and let \(\theta:F_2\to K\) be a
finite collection of simultaneous word evaluations.
If \(D=\theta(F_2')\) is perfect, then
\[
 (\eta,\theta)(F_2')=S'\times D.
 \tag{C.28a}
\]
Thus every prescribed \(s\in S'\) has a commutator-word
representative with \(\theta(f)=1\).

This applies to the five pentagon evaluations and
their two- and three-cycle substitutions in the
following actual braid quotient in
\({\rm PGL}_3(\mathbb F_{11})\):
\[
\begin{split}
 \sigma_1&\longmapsto
 \begin{pmatrix}-1&2&0\\1&0&0\\0&0&1\end{pmatrix},\qquad
 \sigma_2\longmapsto
 \begin{pmatrix}1&0&0\\0&-1&2\\0&1&0\end{pmatrix},\\
 \sigma_3&\longmapsto
 \begin{pmatrix}1&0&-1\\0&1&-1\\0&0&-2\end{pmatrix}.
\end{split}
\tag{C.28b}
\]
Its common derived image is perfect, of order
\(120\cdot1320^3=275996160000\).
Consequently this finite pentagon test excludes no
prescribed conjugator coordinate in any finite solvable
source, including the pure source inside \(G_6\).
This is a simultaneous finite-image statement; it
does not construct a full-GT lift of that coordinate.
\end{statement}
\begin{proof}
The joint image in (C.28a) is a subdirect product.
To recall the elementary subdirect-product argument,
for \(J\leq A\times B\) surjecting on both factors,
the subgroups
\(A_0=\{a:(a,1)\in J\}\) and
\(B_0=\{b:(1,b)\in J\}\) are normal, and
\((a,b)\in J\) identifies \(A/A_0\) with
\(B/B_0\). Apply this to \(A=S'\), \(B=D\).
The common quotient is both solvable and perfect,
so it is trivial. This proves (C.28a), for actual
finite words as well as for their profinite completion.

Here are the complete small certificates establishing
the hypothesis for (C.28b). The matrices are obtained
by quotienting the usual four-dimensional braid
matrices, with parameter \(t=2\), by their fixed
vector \((1,1,1,1)\). Direct matrix multiplication
also independently checks all braid relations and
the distinction of the two endpoint matrices.
Use the five pairs from C.27, followed for each
\((x,y)\) by \((y,x),(z,x),(y,z)\), with
\(z=(xy)^{-1}\). There are twenty indexed coordinates.
The executed certificate partitions them as follows:
\[
\begin{array}{c|l|r|r}
 j&\text{zero-based coordinate indices}&|D_j|&|D_j/Z(D_j)|\\\hline
 0&0,2,5,6,7,11,12,13&120&60\\
 1&1,8,16,18&1320&660\\
 2&3,9,17,19&1320&660\\
 3&4,10,14,15&1320&660
\end{array}
\tag{C.28c}
\]
For each block, normal-closure enumeration from the
single source commutator gives exactly the displayed
number of elements. Its first coordinate distinguishes
all of them, so every other coordinate in the block
is a homomorphic function of that one. The program
finds generators, enumerates their commutator normal
closure, and obtains the entire group again.
It then enumerates the centre, which has order two,
and the conjugacy classes of the central quotient.
The class sizes are
\[
 (1,12,12,15,20),\qquad
 (1,55,60,60,110,110,132,132),
\]
respectively. For every nonidentity class it
independently verifies that a representative's normal
closure is the whole quotient. Thus all four
\(D_j\) are perfect central extensions of the indicated
nonabelian simple groups.

It remains to exclude diagonal identifications between
the three simple quotients of order 660. Write
\(c=[u,v]\) in the source free group. The literal words
\[
 w_u=ucu^{-1}c,\qquad w_v=vcv^{-1}c
\]
provide exact order witnesses. On factors \((1,2)\),
\(w_u\) has quotient orders \((5,11)\); on factors
\((2,3)\) its orders are \((11,5)\); on factors
\((1,3)\), \(w_v\) has orders \((5,11)\).
Therefore none of these pairs is the graph of an
isomorphism. The order-60 simple quotient cannot be
identified with an order-660 quotient.

For completeness, a subdirect subgroup of a product
of nonabelian simple groups whose every two-factor
projection is the full product is the full product.
Induct on the number of factors. Applying the same
subdirect-product argument to the product of the
first factors and the last simple group, a proper
subdirect product would give a simple quotient of
the first product. Any normal subgroup of a product
of nonabelian simple groups is a product of factors:
commuting with each factor shows that each nontrivial
coordinate supplies that entire factor. The simple
quotient would therefore identify the last factor
with one earlier factor, contrary to the two-factor
condition. This proves the claim.

It follows that the actual common derived image
\(D\leq P=\prod_{j=0}^3D_j\) surjects onto
\(P/Z(P)\), so \(DZ(P)=P\). Since \(P\) is perfect,
\[
 P=P'=(DZ(P))'=D'\leq D\leq P.
\]
Hence \(D=P\), proving its order and perfectness.
The independent program
\path{c_burau_perfect_certificate.py}
uses only exact projective matrix arithmetic and
these four small enumerations; its JSON includes
all class representatives, centres, generators, and
the two separating words. The separate permutation
calculation \path{c_burau_pentagon_probe.py} gives
the same common derived order and second derived order.
The option \texttt{--full-orbits} repeats the independent
certificate after including every one of C.27's six
substitutions and all their equation coordinates.
There are then sixty coordinates, partitioned into
blocks of sizes \(24,12,12,12\); their derived groups
and the separating words are unchanged. This complete
repeat was run successfully, so the conclusion also
holds with all those necessary gentle identities imposed.

Finally, the pure source inside \(G_6\) is solvable.
Its kernel over the congruence image is a subgroup
of \(N/N'N^6\), hence abelian, and its congruence
image is a subgroup of the solvable group
\(Q_{576}\). For the latter assertion,
\({\rm SL}_2(\mathbb Z/4)\) is an extension of
\(S_3\) by an elementary abelian 2-group, and
\({\rm SL}_2(\mathbb F_3)\) is an extension of
\(C_3\) by \(Q_8\); quotienting their product by
the common scalar sign preserves solvability.
Apply (C.28a) while keeping the entire prescribed
pure-source value. Setting all sixty braid values
to one satisfies every tested gentle equation and
the pentagon. No inference of compatibility at all
finite levels is made.
\end{proof}

